\section{Introduction}\label{sec:introduction}
Correlations are increasingly recognised as decisive failure modes in quantum error correction (QEC)~\cite{harperLearningCorrelatedNoise2023,fkamDetrimentalNonMarkovianErrors2025}. Even when \emph{average} error rates lie below conventional thresholds, the physics of real devices permits temporal memory and spatial structure that can concentrate faults into bursts, induce heavy tails in syndrome statistics, and ultimately undermine the effective independence assumptions behind standard threshold theorems. At the same time, the modelling tools on either side of the physics--QEC interface remain misaligned. On one hand, independent stochastic Pauli models are the workhorse of QEC design and analysis, enabling scalable simulation, decoding, and the threshold statements that have driven much recent progress. On the other hand, open quantum systems models can capture coherent dynamics and non-Markovian memory, but are rarely presented in a form directly compatible with circuit-level simulation and benchmarking. What is missing is a common language that (i) is operationally tied to the noise seen by stabiliser circuits, (ii) retains genuinely multi-time and multi-qubit correlations, and (iii) admits efficient representations suitable for simulation, inference, and comparison across devices.

A natural starting point for describing general multi-time noise is the quantum \emph{process tensor}~\cite{pollockNonMarkovianQuantumProcesses2018}. Conceptually, it is the multi-time generalisation of a quantum channel. It provides an operationally complete description of non-Markovian dynamics by encoding the joint statistics of the system's response to arbitrary interventions. It thereby captures the temporal correlation structure induced by an environment with memory. However, this generality comes with two drawbacks for QEC. First, process tensors are high-rank quantum objects whose complexity scales exponentially with both time and system dimension. Second, they do not naturally translate into the discrete Pauli fault language used for circuit-level simulation and decoding.

To bridge this gap, we introduce \emph{spatiotemporal Pauli processes} (SPPs): stochastic processes over Pauli trajectories obtained by applying a \emph{multi-time Pauli twirl} to a general process tensor. Operationally, this construction formalises what is induced by standard Pauli-frame randomisation and related randomised compiling techniques. Viewing the multi-time Pauli twirl as a \emph{superprocess}, our central conceptual result is that it maps any (possibly non-Markovian) quantum process to a \emph{process-separable} Pauli comb---equivalently, a well-defined joint probability distribution over Pauli operator trajectories across space and time. In this sense, SPPs are the natural multi-time analogue of Pauli channels. This provides an operationally grounded noise framework compatible with both microscopic open systems modelling and the stabiliser workflow of QEC design and analysis.

\Cref{fig:paper-overview} provides a schematic overview of the framework including the key technical results and applications developed in this work.

\begin{figure}[t]
    \centering
    \begin{adjustbox}{width=\textwidth, keepaspectratio}
        \begin{tikzpicture}
	\begin{pgfonlayer}{nodelayer}
		\node [style=none] (12) at (4.9, 4.3) {};
		\node [style=none] (13) at (4.9, 3.8) {};
		\node [style=none] (14) at (4.9, 3.3) {};
		\node [style=circgate] (15) at (6.05, 4.3) {};
		\node [style=circgate] (20) at (7.65, 3.3) {};
		\node [style=circgate] (21) at (9.25, 4.3) {};
		\node [style=circgate] (22) at (9.25, 3.8) {};
		\node [style=none] (24) at (10.4, 4.3) {};
		\node [style=none] (25) at (10.4, 3.8) {};
		\node [style=none] (26) at (10.4, 3.3) {};
		\node [style=none] (44) at (4.9, 2.8) {};
		\node [style=circgate] (45) at (6.05, 2.8) {};
		\node [style=circgate] (46) at (7.65, 2.8) {};
		\node [style=none] (48) at (10.4, 2.8) {};
		\node [style=circpauli2] (49) at (5.675, 4.3) {};
		\node [style=circpauli0] (50) at (5.675, 3.8) {};
		\node [style=circpauli1] (51) at (6.425, 4.3) {};
		\node [style=circpauli0] (52) at (6.425, 3.8) {};
		\node [style=circpauli3] (53) at (6.425, 3.3) {};
		\node [style=circpauli3] (54) at (5.675, 3.3) {};
		\node [style=circpauli1] (55) at (5.675, 2.8) {};
		\node [style=circpauli0] (56) at (6.425, 2.8) {};
		\node [style=circpauli1] (57) at (7.275, 4.3) {};
		\node [style=circpauli0] (58) at (7.275, 3.8) {};
		\node [style=circpauli2] (59) at (8.025, 4.3) {};
		\node [style=circpauli3] (60) at (8.025, 3.8) {};
		\node [style=circpauli1] (61) at (8.025, 3.3) {};
		\node [style=circpauli3] (62) at (7.275, 3.3) {};
		\node [style=circpauli0] (63) at (7.275, 2.8) {};
		\node [style=circpauli2] (64) at (8.025, 2.8) {};
		\node [style=circpauli2] (65) at (8.875, 4.3) {};
		\node [style=circpauli3] (66) at (8.875, 3.8) {};
		\node [style=circpauli3] (67) at (9.625, 4.3) {};
		\node [style=circpauli1] (68) at (9.625, 3.8) {};
		\node [style=circpauli0] (69) at (9.625, 3.3) {};
		\node [style=circpauli1] (70) at (8.875, 3.3) {};
		\node [style=circpauli2] (71) at (8.875, 2.8) {};
		\node [style=circpauli0] (72) at (9.625, 2.8) {};
		\node [style=none] (27) at (7.075, 2.525) {};
		\node [style=none] (28) at (6.625, 2.525) {};
		\node [style=none] (29) at (5.475, 2.525) {};
		\node [style=none] (30) at (5.025, 2.525) {};
		\node [style=none] (31) at (5.025, 5.075) {};
		\node [style=none] (33) at (6.625, 4.575) {};
		\node [style=none] (34) at (5.475, 4.575) {};
		\node [style=none] (35) at (8.225, 2.525) {};
		\node [style=none] (36) at (8.675, 2.525) {};
		\node [style=none] (37) at (9.825, 2.525) {};
		\node [style=none] (38) at (10.275, 2.525) {};
		\node [style=none] (39) at (10.275, 5.075) {};
		\node [style=none] (40) at (9.825, 4.575) {};
		\node [style=none] (41) at (8.675, 4.575) {};
		\node [style=none] (42) at (8.225, 4.575) {};
		\node [style=none] (43) at (7.075, 4.575) {};
		\node [style=circgatetall] (73) at (6.05, 3.55) {};
		\node [style=circgatetall] (74) at (7.65, 4.05) {};
		\node [style=circgatetall] (75) at (9.25, 3.05) {};
		\node [style=none] (76) at (7.65, 4.825) {\(\Upsilon_{0:k}\)};
		\node [style=none] (77) at (0.325, 4.3) {};
		\node [style=none] (78) at (0.325, 3.8) {};
		\node [style=none] (79) at (0.325, 3.3) {};
		\node [style=circgate] (80) at (1.175, 4.3) {};
		\node [style=circgate] (81) at (2.175, 3.3) {};
		\node [style=circgate] (82) at (3.175, 4.3) {};
		\node [style=circgate] (83) at (3.175, 3.8) {};
		\node [style=none] (84) at (4.025, 4.3) {};
		\node [style=none] (85) at (4.025, 3.8) {};
		\node [style=none] (86) at (4.025, 3.3) {};
		\node [style=none] (87) at (0.325, 2.8) {};
		\node [style=circgate] (88) at (1.175, 2.8) {};
		\node [style=circgate] (89) at (2.175, 2.8) {};
		\node [style=none] (90) at (4.025, 2.8) {};
		\node [style=none] (115) at (1.9, 2.525) {};
		\node [style=none] (116) at (1.45, 2.525) {};
		\node [style=none] (117) at (0.9, 2.525) {};
		\node [style=none] (118) at (0.45, 2.525) {};
		\node [style=none] (119) at (0.45, 5.075) {};
		\node [style=none] (120) at (1.45, 4.575) {};
		\node [style=none] (121) at (0.9, 4.575) {};
		\node [style=none] (122) at (2.45, 2.525) {};
		\node [style=none] (123) at (2.9, 2.525) {};
		\node [style=none] (124) at (3.45, 2.525) {};
		\node [style=none] (125) at (3.9, 2.525) {};
		\node [style=none] (126) at (3.9, 5.075) {};
		\node [style=none] (127) at (3.45, 4.575) {};
		\node [style=none] (128) at (2.9, 4.575) {};
		\node [style=none] (129) at (2.45, 4.575) {};
		\node [style=none] (130) at (1.9, 4.575) {};
		\node [style=circgatetall] (131) at (1.175, 3.55) {};
		\node [style=circgatetall] (132) at (2.175, 4.05) {};
		\node [style=circgatetall] (133) at (3.175, 3.05) {};
		\node [style=none] (134) at (2.175, 4.825) {\(\Upsilon_{0:k}\)};
		\node [style=none] (135) at (11.225, 4.3) {};
		\node [style=none] (136) at (11.225, 3.8) {};
		\node [style=none] (137) at (11.225, 3.3) {};
		\node [style=circgate] (138) at (12.075, 4.3) {};
		\node [style=circgate] (139) at (13.075, 3.3) {};
		\node [style=circgate] (140) at (14.075, 4.3) {};
		\node [style=circgate] (141) at (14.075, 3.8) {};
		\node [style=none] (142) at (14.925, 4.3) {};
		\node [style=none] (143) at (14.925, 3.8) {};
		\node [style=none] (144) at (14.925, 3.3) {};
		\node [style=none] (145) at (11.225, 2.8) {};
		\node [style=circgate] (146) at (12.075, 2.8) {};
		\node [style=circgate] (147) at (13.075, 2.8) {};
		\node [style=none] (148) at (14.925, 2.8) {};
		\node [style=circgatetall] (165) at (12.075, 3.55) {};
		\node [style=circgatetall] (166) at (13.075, 4.05) {};
		\node [style=circgatetall] (167) at (14.075, 3.05) {};
		\node [style=circpauli4] (169) at (12.575, 3.3) {};
		\node [style=circpauli4] (170) at (12.575, 2.8) {};
		\node [style=circpauli4] (171) at (12.575, 4.3) {};
		\node [style=circpauli4] (172) at (12.575, 3.8) {};
		\node [style=circpauli4] (173) at (11.575, 3.3) {};
		\node [style=circpauli4] (174) at (11.575, 2.8) {};
		\node [style=circpauli4] (175) at (11.575, 4.3) {};
		\node [style=circpauli4] (176) at (11.575, 3.8) {};
		\node [style=circpauli4] (177) at (14.575, 3.3) {};
		\node [style=circpauli4] (178) at (14.575, 2.8) {};
		\node [style=circpauli4] (179) at (14.575, 4.3) {};
		\node [style=circpauli4] (180) at (14.575, 3.8) {};
		\node [style=circpauli4] (181) at (13.575, 3.3) {};
		\node [style=circpauli4] (182) at (13.575, 2.8) {};
		\node [style=circpauli4] (183) at (13.575, 4.3) {};
		\node [style=circpauli4] (184) at (13.575, 3.8) {};
		\node [style=none] (191) at (12.8, 2.525) {};
		\node [style=none] (192) at (12.35, 2.525) {};
		\node [style=none] (193) at (11.8, 2.525) {};
		\node [style=none] (194) at (11.35, 2.525) {};
		\node [style=none] (195) at (11.35, 5.075) {};
		\node [style=none] (196) at (12.35, 4.575) {};
		\node [style=none] (197) at (11.8, 4.575) {};
		\node [style=none] (198) at (13.35, 2.525) {};
		\node [style=none] (199) at (13.8, 2.525) {};
		\node [style=none] (200) at (14.35, 2.525) {};
		\node [style=none] (201) at (14.8, 2.525) {};
		\node [style=none] (202) at (14.8, 5.075) {};
		\node [style=none] (203) at (14.35, 4.575) {};
		\node [style=none] (204) at (13.8, 4.575) {};
		\node [style=none] (205) at (13.35, 4.575) {};
		\node [style=none] (206) at (12.8, 4.575) {};
		\node [style=footnotesize] (207) at (13.075, 4.825) {\(\mathcal{T}_P^{(k)}[\Upsilon_{0:k}]\)};
		\node [style=none] (208) at (11.575, 4.825) {};
		\node [style=none] (209) at (12.575, 4.825) {};
		\node [style=none] (210) at (13.825, 4.825) {};
		\node [style=none] (211) at (14.575, 4.825) {};
		\node [style=none] (212) at (12.575, 4.55) {};
		\node [style=none] (213) at (13.575, 4.55) {};
		\node [style=none] (215) at (14.2, 4.825) {};
		\node [style=none] (216) at (11.95, 4.825) {};
		\node [style=none] (217) at (12.325, 4.825) {};
		\node [style=none] (218) at (4.2, 3.775) {};
		\node [style=none] (219) at (4.75, 3.775) {};
		\node [style=none] (220) at (10.575, 3.775) {};
		\node [style=none] (221) at (11.1, 3.775) {};
		\node [style=pttensor] (222) at (0.675, 1.275) {};
		\node [style=none] (223) at (0.425, 0.825) {};
		\node [style=none] (224) at (0.925, 0.825) {};
		\node [style=smallsize] (234) at (0.425, 0.7) {\({}_{\alpha_0}\)};
		\node [style=smallsize] (235) at (0.925, 0.7) {\({}_{\beta_0}\)};
		\node [style=pttensor] (236) at (1.675, 1.275) {};
		\node [style=none] (237) at (1.425, 0.825) {};
		\node [style=none] (238) at (1.925, 0.825) {};
		\node [style=smallsize] (239) at (1.425, 0.7) {\({}_{\alpha_1}\)};
		\node [style=smallsize] (240) at (1.925, 0.7) {\({}_{\beta_1}\)};
		\node [style=pttensor] (241) at (2.675, 1.275) {};
		\node [style=none] (242) at (2.425, 0.825) {};
		\node [style=none] (243) at (2.925, 0.825) {};
		\node [style=smallsize] (244) at (2.425, 0.7) {\({}_{\alpha_2}\)};
		\node [style=smallsize] (245) at (2.925, 0.7) {\({}_{\beta_2}\)};
		\node [style=pttensor] (246) at (3.675, 1.275) {};
		\node [style=none] (247) at (3.425, 0.825) {};
		\node [style=none] (248) at (3.925, 0.825) {};
		\node [style=smallsize] (249) at (3.425, 0.7) {\({}_{\alpha_k}\)};
		\node [style=smallsize] (250) at (3.925, 0.7) {\({}_{\beta_k}\)};
		\node [style=none] (251) at (3.35, 1.275) {};
		\node [style=none] (252) at (3, 1.275) {};
		\node [style=pttensorlarge] (253) at (6.775, 1.475) {};
		\node [style=twirltensor] (254) at (6.525, 0.875) {\(P\)};
		\node [style=twirltensor] (255) at (7.025, 0.875) {\(P\)};
		\node [style=none] (256) at (6.525, 0.525) {};
		\node [style=none] (257) at (7.025, 0.525) {};
		\node [style=none] (258) at (7.275, 1.475) {};
		\node [style=none] (259) at (6.275, 1.475) {};
		\node [style=pttensorlarge] (260) at (8.525, 1.475) {};
		\node [style=twirltensor] (261) at (8.525, 0.875) {\(P\)};
		\node [style=none] (263) at (8.525, 0.525) {};
		\node [style=none] (265) at (9.025, 1.475) {};
		\node [style=none] (266) at (8.025, 1.475) {};
		\node [style=spptensor] (267) at (11.575, 1.275) {};
		\node [style=none] (268) at (11.575, 0.825) {};
		\node [style=smallsize] (270) at (11.575, 0.7) {\({}_{x_0}\)};
		\node [style=spptensor] (272) at (12.575, 1.275) {};
		\node [style=none] (273) at (12.575, 0.825) {};
		\node [style=smallsize] (275) at (12.575, 0.7) {\({}_{x_1}\)};
		\node [style=spptensor] (277) at (13.575, 1.275) {};
		\node [style=none] (278) at (13.575, 0.825) {};
		\node [style=smallsize] (280) at (13.575, 0.7) {\({}_{x_2}\)};
		\node [style=spptensor] (282) at (14.575, 1.275) {};
		\node [style=none] (283) at (14.575, 0.825) {};
		\node [style=smallsize] (285) at (14.575, 0.7) {\({}_{x_k}\)};
		\node [style=none] (287) at (14.25, 1.275) {};
		\node [style=none] (288) at (13.9, 1.275) {};
		\node [style=footnotesize] (289) at (2.175, 5.425) {Process tensor noise model};
		\node [style=footnotesize] (290) at (7.65, 5.425) {Multi-time Pauli twirl};
		\node [style=footnotesize] (291) at (13.075, 5.425) {Spatiotemporal Pauli process};
		\node [style=footnotesize] (292) at (2.175, 2.05) {Process tensor MPO};
		\node [style=footnotesize] (293) at (7.65, 2.05) {Local Pauli twirl contraction};
		\node [style=footnotesize] (294) at (13.075, 2.05) {SPP MPS};
		\node [style=none] (295) at (4.375, 1.15) {};
		\node [style=none] (296) at (5.8, 1.15) {};
		\node [style=none] (297) at (9.525, 1.15) {};
		\node [style=none] (298) at (10.95, 1.15) {};
		\node [style=none] (299) at (15.375, 6.45) {};
		\node [style=none] (300) at (0, 6.45) {};
		\node [style=none] (301) at (15.375, 0.25) {};
		\node [style=none] (302) at (0, 0.25) {};
		\node [style=none] (303) at (7.65, 6) {\textbf{Bridge: Process tensors to spatiotemporal Pauli processes}};
		\node [style=none] (306) at (6.3, -0.3) {};
		\node [style=none] (307) at (0, -0.3) {};
		\node [style=none] (308) at (6.3, -5.25) {};
		\node [style=none] (309) at (0, -5.25) {};
		\node [style=none] (310) at (15.375, -0.3) {};
		\node [style=none] (311) at (6.575, -0.3) {};
		\node [style=none] (312) at (15.375, -5.25) {};
		\node [style=none] (313) at (6.575, -5.25) {};
		\node [style=none] (314) at (3.15, -0.75) {\textbf{Open quantum systems}};
		\node [style=none] (315) at (10.975, -0.75) {\textbf{Analysis and simulation}};
		\node [style=none] (318) at (2.925, -0.25) {};
		\node [style=none] (319) at (3.15, 0.2) {};
		\node [style=none] (320) at (3.375, -0.25) {};
		\node [style=none] (321) at (3.5, -0.025) {};
		\node [style=none] (322) at (2.8, -0.025) {};
		\node [style=none] (323) at (2.925, -0.025) {};
		\node [style=none] (324) at (3.375, -0.025) {};
		\node [style=none] (325) at (10.75, 0.2) {};
		\node [style=none] (326) at (10.975, -0.25) {};
		\node [style=none] (327) at (11.2, 0.2) {};
		\node [style=none] (328) at (11.325, -0.025) {};
		\node [style=none] (329) at (10.625, -0.025) {};
		\node [style=none] (330) at (10.75, -0.025) {};
		\node [style=none] (331) at (11.2, -0.025) {};
		\node [style=none] (333) at (0.5, -3.5) {};
		\node [style=none] (334) at (0.5, -4) {};
		\node [style=circgate] (336) at (1.425, -4) {};
		\node [style=circgate] (338) at (2, -3.5) {};
		\node [style=none] (340) at (2.9, -3.5) {};
		\node [style=none] (341) at (2.9, -4) {};
		\node [style=none] (342) at (0.5, -4.5) {};
		\node [style=circgate] (343) at (0.85, -4.5) {};
		\node [style=circgate] (344) at (1.425, -4.5) {};
		\node [style=none] (345) at (2.9, -4.5) {};
		\node [style=circgate] (346) at (0.85, -3.5) {};
		\node [style=circgate] (349) at (1.425, -3.5) {};
		\node [style=circgate] (350) at (2, -4) {};
		\node [style=node] (353) at (2, -4.5) {};
		\node [style=node] (354) at (0.85, -4) {};
		\node [style=footnotesize] (356) at (1.7, -2.875) {QEC circuit};
		\node [style=footnotesize] (357) at (3.15, -1.4) {System-environment dynamics};
		\node [style=largesize] (358) at (3.15, -2.025) {\(\mathcal{H} = \mathcal{H}_S + \mathcal{H}_E + \mathcal{H}_{SE}\)};
		\node [style=circgate] (359) at (2.575, -4) {};
		\node [style=circgate] (360) at (2.575, -4.5) {};
		\node [style=tallboxyellow] (361) at (5.025, -4) {\(\mathcal{E}\)};
		\node [style=lefttalltriangle] (363) at (4.025, -4) {\(\rho_{{\!}_{S\!E}}\)};
		\node [style=none] (364) at (4.025, -3.675) {};
		\node [style=none] (365) at (4.025, -4.325) {};
		\node [style=none] (368) at (5.725, -3.675) {};
		\node [style=none] (369) at (5.725, -4.325) {};
		\node [style=footnotesize] (370) at (4.65, -2.875) {Quantum channel};
		\node [style=node] (371) at (2.575, -3.5) {};
		\node [style=tinycircleorange] (372) at (7.625, -2.6) {};
		\node [style=tinycircleorange] (373) at (8.125, -2.6) {};
		\node [style=tinycircleorange] (374) at (8.625, -2.6) {};
		\node [style=tinycircleorange] (375) at (9.125, -2.6) {};
		\node [style=tinycircleorange] (376) at (7.95, -2.225) {};
		\node [style=tinycircleorange] (377) at (8.45, -2.225) {};
		\node [style=tinycircleorange] (378) at (8.95, -2.225) {};
		\node [style=tinycircleorange] (379) at (9.45, -2.225) {};
		\node [style=tinycircleorange] (380) at (8.275, -1.85) {};
		\node [style=tinycircleorange] (381) at (8.775, -1.85) {};
		\node [style=tinycircleorange] (382) at (9.275, -1.85) {};
		\node [style=tinycircleorange] (383) at (9.775, -1.85) {};
		\node [style=none] (384) at (7.625, -2.9) {};
		\node [style=none] (385) at (8.125, -2.9) {};
		\node [style=none] (386) at (8.625, -2.9) {};
		\node [style=none] (387) at (9.125, -2.9) {};
		\node [style=none] (388) at (7.95, -2.525) {};
		\node [style=none] (389) at (8.45, -2.525) {};
		\node [style=none] (390) at (8.95, -2.525) {};
		\node [style=none] (391) at (9.45, -2.525) {};
		\node [style=none] (392) at (8.275, -2.15) {};
		\node [style=none] (393) at (8.775, -2.15) {};
		\node [style=none] (394) at (9.275, -2.15) {};
		\node [style=none] (395) at (9.775, -2.15) {};
		\node [style=footnotesize] (396) at (8.7, -1.3) {SPP PEPS};
		\node [style=boxyellowsmall] (397) at (11.425, -2.05) {$T$};
		\node [style=none] (398) at (12.05, -2.05) {};
		\node [style=none] (399) at (10.8, -2.05) {};
		\node [style=spptensorsmall] (400) at (14.225, -1.85) {\(\mathbf{A}\)};
		\node [style=none] (401) at (14.85, -1.85) {};
		\node [style=none] (402) at (13.6, -1.85) {};
		\node [style=none] (403) at (14.225, -2.35) {};
		\node [style=none] (404) at (14.225, -2.475) {\({}_x\)};
		\node [style=largesize] (405) at (12.825, -2.1) {\(=\,\displaystyle\sum_x\)};
		\node [style=footnotesize] (406) at (12.8, -1.3) {Transfer operator};
		\node [style=none] (407) at (12.8, -4.175) {\includegraphics{figures/experiment4_mini.pdf}};
		\node [style=footnotesize] (408) at (12.8, -2.85) {QEC simulation};
		\node [style=hmmred] (409) at (9.3, -3.875) {\(\mathbf{1}\)};
		\node [style=hmmblue] (410) at (8.275, -4.75) {\(\mathbf{0}\)};
		\node [style=footnotesize] (411) at (8.775, -3.3) {Hidden Markov model};
		\node [style=none] (412) at (9.25, -4.65) {\({}_{\Gamma_{01}}\)};
		\node [style=none] (413) at (8.5, -3.9) {\({}_{\Gamma_{10}}\)};
		\node [style=none] (414) at (10.25, -4.075) {\({}_{\Gamma_{11}}\)};
		\node [style=none] (415) at (7.375, -4.575) {\({}_{\Gamma_{00}}\)};
	\end{pgfonlayer}
	\begin{pgfonlayer}{edgelayer}
		\draw (12.center) to (24.center);
		\draw (25.center) to (13.center);
		\draw (14.center) to (26.center);
		\draw (44.center) to (48.center);
		\draw [style=drawfillblue] (29.center)
			 to (30.center)
			 to (31.center)
			 to (39.center)
			 to (38.center)
			 to (37.center)
			 to (40.center)
			 to (41.center)
			 to (36.center)
			 to (35.center)
			 to (42.center)
			 to (43.center)
			 to (27.center)
			 to (28.center)
			 to (33.center)
			 to (34.center)
			 to cycle;
		\draw (77.center) to (84.center);
		\draw (85.center) to (78.center);
		\draw (79.center) to (86.center);
		\draw (87.center) to (90.center);
		\draw [style=drawfillblue] (117.center)
			 to (118.center)
			 to (119.center)
			 to (126.center)
			 to (125.center)
			 to (124.center)
			 to (127.center)
			 to (128.center)
			 to (123.center)
			 to (122.center)
			 to (129.center)
			 to (130.center)
			 to (115.center)
			 to (116.center)
			 to (120.center)
			 to (121.center)
			 to cycle;
		\draw (135.center) to (142.center);
		\draw (143.center) to (136.center);
		\draw (137.center) to (144.center);
		\draw (145.center) to (148.center);
		\draw [style=thickorange] (192.center)
			 to (196.center)
			 to (197.center)
			 to (193.center)
			 to (194.center)
			 to (195.center)
			 to (202.center)
			 to (201.center)
			 to (200.center)
			 to (203.center)
			 to (204.center)
			 to (199.center)
			 to (198.center)
			 to (205.center)
			 to (206.center)
			 to (191.center)
			 to cycle;
		\draw [style=orangewire] (178)
			 to (211.center)
			 to (215.center);
		\draw [style=orangewire] (174)
			 to (208.center)
			 to (216.center);
		\draw [style=orangedotted] (183) to (213.center);
		\draw [style=orangedotted] (171) to (212.center);
		\draw [style=orangedotted] (215.center) to (210.center);
		\draw [style=orangedotted] (216.center) to (217.center);
		\draw [style=orangewire] (182) to (213.center);
		\draw [style=orangewire] (170) to (212.center);
		\draw [style=arrow] (218.center) to (219.center);
		\draw [style=arrow] (220.center) to (221.center);
		\draw [in=90, out=-45] (222) to (224.center);
		\draw [in=90, out=-135] (222) to (223.center);
		\draw [in=90, out=-45] (236) to (238.center);
		\draw [in=90, out=-135] (236) to (237.center);
		\draw [in=90, out=-45] (241) to (243.center);
		\draw [in=90, out=-135] (241) to (242.center);
		\draw [in=90, out=-45] (246) to (248.center);
		\draw [in=90, out=-135] (246) to (247.center);
		\draw [style=thickwire] (222) to (236);
		\draw [style=thickwire] (236) to (241);
		\draw [style=thickwire] (252.center) to (241);
		\draw [style=thickwire] (251.center) to (246);
		\draw [style=thickdotted] (252.center) to (251.center);
		\draw [in=90, out=-45] (253) to (255);
		\draw [in=90, out=-135] (253) to (254);
		\draw (254) to (255);
		\draw (257.center) to (255);
		\draw (256.center) to (254);
		\draw [style=thickwire] (253) to (258.center);
		\draw [style=thickwire] (253) to (259.center);
		\draw (263.center) to (261);
		\draw [style=thickwire] (260) to (265.center);
		\draw [style=thickwire] (260) to (266.center);
		\draw [bend left] (261) to (260);
		\draw [bend right] (261) to (260);
		\draw [style=thickwire] (267) to (272);
		\draw [style=thickwire] (272) to (277);
		\draw [style=thickwire] (288.center) to (277);
		\draw [style=thickwire] (287.center) to (282);
		\draw [style=thickdotted] (288.center) to (287.center);
		\draw (282) to (283.center);
		\draw (278.center) to (277);
		\draw (273.center) to (272);
		\draw (268.center) to (267);
		\draw [style=arrow] (295.center) to (296.center);
		\draw [style=arrow] (297.center) to (298.center);
		\draw [style=thickgreen] (302.center)
			 to (301.center)
			 to (299.center)
			 to (300.center)
			 to cycle;
		\draw [style=thickred] (309.center)
			 to (308.center)
			 to (306.center)
			 to (307.center)
			 to cycle;
		\draw [style=thickblue] (312.center)
			 to (310.center)
			 to (311.center)
			 to (313.center)
			 to cycle;
		\draw [style=fillredoutline] (320.center)
			 to (318.center)
			 to (323.center)
			 to (322.center)
			 to (319.center)
			 to (321.center)
			 to (324.center)
			 to cycle;
		\draw [style=fillblueoutline] (327.center)
			 to (325.center)
			 to (330.center)
			 to (329.center)
			 to (326.center)
			 to (328.center)
			 to (331.center)
			 to cycle;
		\draw (340.center) to (333.center);
		\draw (334.center) to (341.center);
		\draw (342.center) to (345.center);
		\draw (353) to (350);
		\draw (354) to (343);
		\draw (364.center) to (368.center);
		\draw (365.center) to (369.center);
		\draw (359) to (371);
		\draw (380) to (383);
		\draw (376) to (379);
		\draw (372) to (375);
		\draw (380) to (372);
		\draw (381) to (373);
		\draw (382) to (374);
		\draw (383) to (375);
		\draw (384.center) to (372);
		\draw (388.center) to (376);
		\draw (392.center) to (380);
		\draw (385.center) to (373);
		\draw (389.center) to (377);
		\draw (393.center) to (381);
		\draw (386.center) to (374);
		\draw (390.center) to (378);
		\draw (394.center) to (382);
		\draw (387.center) to (375);
		\draw (391.center) to (379);
		\draw (395.center) to (383);
		\draw [style=fusedwire] (397) to (398.center);
		\draw [style=fusedwire] (399.center) to (397);
		\draw [style=fusedwire] (400) to (401.center);
		\draw [style=fusedwire] (402.center) to (400);
		\draw (403.center) to (400);
		\draw [style=arrow, bend right] (409) to (410);
		\draw [style=arrow, bend right] (410) to (409);
		\draw [style=arrow, in=-60, out=0, loop] (409) to ();
		\draw [style=arrow, in=120, out=-180, loop] (410) to ();
	\end{pgfonlayer}
\end{tikzpicture}

    \end{adjustbox}
    \caption{\textbf{Framework overview: Process tensors to spatiotemporal Pauli processes.}}
    \label{fig:paper-overview}
\end{figure}

Beyond the operational definition, we show that SPPs admit efficient tensor network representations that can be constructed directly from an underlying system-environment interaction. Because the multi-time Pauli twirl acts strictly via local contractions on the system legs, the induced SPP inherits a classical tensor network form---whose entries encode joint spatiotemporal Pauli trajectory probabilities. Moreover, the inner bond dimension of this network is bounded by the Liouville-space dimension of the underlying environment, tying the effective memory of the Pauli process to the size of the physical bath. To make these correlations interpretable, we develop a transfer operator formalism for SPPs that links spectral properties to temporal correlation lengths, and demonstrate for suitable classes of SPPs, these tensor networks can be recast exactly to classical hidden Markov models. Together, these structures support efficient Monte Carlo sampling and inference and provide a direct interface between microscopic non-Markovian dynamics and circuit-level QEC simulation.

We deploy this framework to systematically benchmark primitives of surface code computation under correlations. First, using a one-dimensional temporal `storm' model, we demonstrate that tuning the correlation length---while fixing single-round marginal error rates---progressively degrades logical memory and stability, blunting the effective exponential suppression of errors expected from distance scaling. Second, we introduce a genuinely spatiotemporal SPP derived from a two-dimensional coherent quantum cellular automaton (QCA) bath. Under operational Pauli twirling, the QCA dynamics maps to an effective SPP described exactly by a probabilistic cellular automaton (PCA) with a nonlinear transition kernel. Strikingly, tuning a single parameter controlling coherent bath interactions drives the Pauli process into a pseudo-critical regime. In this window, macroscopic error avalanches and critical slowing down cause a complete reversal and breakdown of standard surface code distance scaling. These results establish a concrete, computationally tractable link between non-equilibrium statistical mechanics and QEC performance under spatiotemporally correlated noise.

(\Fidel{ITEMIZED CONTRIBUTIONS + POSITIONING RELATIVE TO OTHER WORK.})

The remainder of this paper is structured as follows. \Cref{sec:background} reviews the necessary background on quantum channels, Pauli twirling, and process tensors. In \cref{sec:pt-tn} we detail the construction of tensor network representations for multi-time processes directly from microscopic dynamics. \Cref{sec:spp} formally introduces SPPs via the multi-time Pauli twirl and derives their tensor network structure. We also illustrate the construction with worked examples. \Cref{sec:correlation-structure} develops the transfer operator formalism and establishes the connection to hidden Markov models. \Cref{sec:temporalstormmodel} we apply this framework to benchmark surface code memory and stability under the temporally correlated storm model. \Cref{sec:qcamodel} extends this to fully spatiotemporal noise, deriving an effective PCA from a coherent quantum bath and demonstrating the resulting breakdown of distance scaling near pseudo-criticality. Finally, \cref{sec:discussion} concludes with a discussion of future applications, including correlation-aware decoding and noise characterisation.

